Optical lattice and single-ion clocks now reach fractional systematic
uncertainties at or below \selfEvalBound---including
ytterbium~\citep{mcgrew2018geodesy, zhu2026yb} and
strontium~\citep{aeppli2024sr, jia2026sr, lu2025ntsc} lattice clocks
and Al$^+$~\citep{marshall2025alplus, dawel2026ptbalplus},
Lu$^+$~\citep{arnold2026lu} and Ca$^+$~\citep{zhang2026ca} ion
clocks---two orders of magnitude better
than current primary frequency standards~\citep{heavner2014nistf2,
weyers2018ptbcsf}. By
comparing independent optical clocks located in different regions and
laboratories, the precision budget self-assessed by individual clocks
can be verified more objectively. This is because the environmental
conditions and geodetic assumptions used in self-evaluation are
typically based on a single location, and using identical methods within
a single laboratory may lead to common systematic biases. Therefore,
connecting clocks from different locations and types into a
collaborative network is a key pathway toward establishing optical
clocks as a verifiable metrological
infrastructure~\citep{riehle2017networks, komar2014quantumnetwork}. The
idea has moved from proposal to practice: fiber links with noise levels
below that of the clocks themselves now span thousands of
kilometers~\citep{schioppo2022cavities, chen2026dissemination}, and
multi-clock campaigns have measured tens of ratios among clocks in six
countries~\citep{lindvall2025coordinated} and across a continental
network~\citep{pizzocaro2026european}. Satellite-based advanced links
with comparable performance are also being developed and will
interconnect with fiber networks to form a future global optical clock
network~\citep{shen2021satellite, caldwell2023geo, shen2022freespace}.

The most immediate purpose of such a network is the redefinition of the
SI second. One of the mandatory requirements (I.2.b) in the roadmap for
this redefinition~\citep{dimarcq2024roadmap} is that frequency ratios of
different clock types be measured at least twice by different institutes
in agreement with an overall uncertainty of the comparison
$\Delta\nu/\nu < \refThreshold$---a requirement reiterated in the 2025
CCTF progress report~\citep{cctf2025strategy, cctf2025taskforce}.
However, comparisons across independent institutions and different
locations have reached at best a precision of
\refEuBestU~\citep{pizzocaro2026european}. So far, no frequency-ratio
measurements between clocks in different cities have satisfied this
requirement~\citep{fortier2026review}.

Closing that gap matters beyond the definition itself. Chronometrically
leveled clocks map the gravitational potential of the Earth at
centimeter resolution~\citep{grotti2018geodesy, takamoto2020redshift,
mcgrew2018geodesy}; clock networks have been proposed and used as
sensors for dark matter~\citep{derevianko2014topological,
wcislo2018global}. All of these applications assume that remote clocks
agree at a level that, until now, has only been demonstrated locally
twice: a Sr/Yb/Al$^+$ triple-clock ratio measurement with an uncertainty
of \refSameCampusU~by NIST--JILA in Boulder (US)~\citep{aeppli2026nist},
and a Yb$^+$/Sr measurement with \refPtbTransportableU~by PTB in
Braunschweig (Germany)~\citep{vishwakarma2026interspeciesclockcomparison5}.

Here we report an intercity, cross-species frequency-ratio measurement
between an IAPMST ytterbium lattice clock in Wuhan and a USTC strontium
lattice clock in Shanghai, about $\siteKm$\,km apart, connected by a
$\fibreKm$\,km fiber link. We determine the Yb/Sr ratio with a total
uncertainty of \uTotalShort$\times10^{-18}$; gravitational redshift
effects have been corrected, whereas measurement error and the
time-varying gravitational redshift term associated with gravitational
tides have been retained as uncertainty components. The result
establishes the two sites as a prototype node pair for a distributed
optical-clock network, and completes, at intercity scale, the class of
ratio measurements that the redefinition process depends on.
